\section{Iterative Inorder Traversal}
\label{sec:trees-iter-inorder}

\problemstatement{Return the inorder traversal (left, node, right) of a
binary tree without recursion, using an explicit stack.}

\begin{patternbox}
Inorder's record step is \emph{delayed}: a node cannot be output until its
entire left subtree is done. So the iterative version has two phases that
alternate -- \textbf{dive left pushing everything}, then \textbf{pop,
record, and turn right}. This ``go as far left as possible, then process,
then step right'' loop is worth memorising; it reappears in the BST
iterator.
\end{patternbox}

\subsection*{Understanding the problem carefully}

In inorder we must reach the leftmost node before recording anything. The
stack stores the chain of ancestors we still owe a visit to. We push
nodes as we descend left; when we can descend no further, the top of the
stack is the next node to record; after recording it we move to its right
subtree and repeat the leftward dive from there.

\begin{ideabox}[A Pointer and a Stack, Cooperating]
Keep a running pointer \textit{curr}. While \textit{curr} is non-null,
push it and move \textit{curr} to its left child (the leftward dive). When
\textit{curr} becomes null, pop a node, record it, and set \textit{curr}
to that node's right child (turn right, then dive left again). The loop
ends only when both the stack is empty \emph{and} \textit{curr} is null.
\end{ideabox}

\subsection*{Algorithm}

\begin{enumerate}
  \item Set \textit{curr} $=$ root; use an empty stack.
  \item While \textit{curr} is non-null \textbf{or} the stack is
        non-empty:
  \begin{enumerate}
    \item While \textit{curr} is non-null: push \textit{curr}, set
          \textit{curr} $=$ \textit{curr}.left.
    \item Pop a node; append its value; set \textit{curr} $=$
          node.right.
  \end{enumerate}
\end{enumerate}

\subsection*{Dry run}

\dryrun{Tree: root $2$; left child $1$; right child $3$.}

\begin{itemize}
  \item \textit{curr}$=2$: push $2$, \textit{curr}$=1$. Push $1$,
        \textit{curr}$=$null.
  \item Pop $1$, record $1$, \textit{curr}$=$null (right of $1$).
  \item Pop $2$, record $2$, \textit{curr}$=3$.
  \item Push $3$, \textit{curr}$=$null. Pop $3$, record $3$.
  \item Result: $1,2,3$.
\end{itemize}

\subsection*{C++ implementation}

\begin{lstlisting}
#include <bits/stdc++.h>
using namespace std;

struct TreeNode {
    int val;
    TreeNode* left;
    TreeNode* right;
    TreeNode(int v) : val(v), left(nullptr), right(nullptr) {}
};

vector<int> inorderIterative(TreeNode* root) {
    vector<int> out;
    stack<TreeNode*> st;
    TreeNode* curr = root;

    while (curr != nullptr || !st.empty()) {
        while (curr != nullptr) {   // dive left
            st.push(curr);
            curr = curr->left;
        }
        curr = st.top();            // leftmost pending node
        st.pop();
        out.push_back(curr->val);   // record
        curr = curr->right;         // turn right
    }
    return out;
}

int main() {
    TreeNode* root = new TreeNode(2);
    root->left = new TreeNode(1);
    root->right = new TreeNode(3);

    for (int x : inorderIterative(root)) cout << x << " ";
    cout << "\n";
    return 0;
}
\end{lstlisting}

\begin{complexitybox}
\textbf{Time Complexity.} $O(n)$ -- each node is pushed once and popped
once. \textbf{Space Complexity.} $O(h)$ -- the stack holds at most one
node per level along the current leftmost chain.
\end{complexitybox}

\begin{takeawaybox}
The ``dive left, pop-record, turn right'' loop is the canonical iterative
inorder and the exact engine behind a BST iterator that returns sorted
values one at a time. Memorise the two-part loop shape; the same structure
handles ``next smallest'' queries on demand.
\end{takeawaybox}
